\documentclass{article}
\usepackage[utf8]{inputenc}
\usepackage[spanish]{babel}
\usepackage{listings}
\usepackage{xcolor}
\usepackage{geometry}
\geometry{a4paper, margin=1in}

\% Configuración para código C++
\definecolor{codegreen}{rgb}{0,0.6,0}        
\definecolor{codegray}{rgb}{0.5,0.5,0.5}     
\definecolor{codepurple}{rgb}{0.58,0,0.82}   
\definecolor{backcolour}{rgb}{0.95,0.95,0.92}

\lstset{
    backgroundcolor=\color{backcolour},      
    commentstyle=\color{codegreen},
    keywordstyle=\color{magenta},
    numberstyle=\tiny\color{codegray},       
    stringstyle=\color{codepurple},
    basicstyle=\ttfamily\small,
    breakatwhitespace=false,
    breaklines=true,
    captionpos=b,
    keepspaces=true,
    numbers=left,
    numbersep=5pt,
    showspaces=false,
    showstringspaces=false,
    showtabs=false,
    tabsize=2,
    language=C++,
    inputencoding=utf8,
    extendedchars=true,
    literate={á}{{\'a}}1 {é}{{\'e}}1 {í}{{\'i}}1 {ó}{{\'o}}1 {ú}{{\'u}}1 {ñ}{{\~n}}1 {¿}{{\textquestiondown}}1 {¡}{{\textexclamdown}}1
}

\% Macros para las secciones de los apuntes
\newcommand{\quees}[1]{
    \section*{🤔 ¿QUÉ ES?}
    \hrulefill \\
    \#1
}

\newcommand{\dichodeotraforma}[1]{
    \section*{💬 DICHO DE OTRA FORMA}
    \hrulefill \\
    \#1
}

\newcommand{\diagrama}[1]{
    \section*{🎨 DIAGRAMA}
    \hrulefill \\
    \#1
}

\newcommand{\comofunciona}[1]{
    \section*{⚙️ ¿CÓMO FUNCIONA?}
    \hrulefill \\
    \#1
}

\newcommand{\complejidad}[1]{
    \section*{📱 COMPLEJIDAD (Big-O)}
    \hrulefill \\
    \#1
}

\newcommand{\entrevistas}[1]{
    \section*{🎯 EN ENTREVISTAS}
    \hrulefill \\
    \#1
}

\newcommand{\resumen}[1]{
    \section*{📋 RESUMEN RÁPIDO}
    \hrulefill \\
    \#1
}

\begin{document}

\title{Stacks (Pilas)}
\author{Aldo Zepeda Montoya}
\date{\today}
\maketitle

\subsection*{CATEGORÍA: 01\_Estructuras\_Lineales}

\quees{
Una pila de platos 🍽️ en el fregadero. Solo puedes poner platos
encima (push) y solo puedes agarrar el de hasta arriba (pop). 

Si quieres el plato de abajo, debes quitar todos los de encima primero. 
Es una estructura "LIFO" (Last In, First Out): el último en entrar
es el primero en salir.

Imagina un tubo de papas Pringles 🥔: la última papa que metieron
en la fábrica es la primera que te comes al abrir el tubo.
}

\dichodeotraforma{
Un Stack es una estructura de datos lineal de tipo ADT (Abstract Data
Type) que restringe el acceso a un solo extremo: el "Top". 

Es fundamental para el funcionamiento de las computadoras, ya que
el "Call Stack" de los lenguajes de programación usa esta lógica
para gestionar las llamadas a funciones y la recursión.
}

\diagrama{\begin{verbatim}
push(A) → push(B) → push(C)

     [C]  ← TOP (C es el último en entrar)
     [B]
     [A]  ← BOTTOM

pop() → devuelve C y lo quita, queda:

     [B]  ← TOP (ahora B es el nuevo tope)
     [A]
\end{verbatim}}

\begin{lstlisting}[language=C++]
1. push(x): Agregas el elemento `x` al tope de la pila. O(1).
2. pop(): Quitas el elemento del tope. O(1).
3. peek() / top(): Miras qué hay en el tope sin quitarlo. O(1).
4. isEmpty(): Verificas si la pila está vacía. O(1).
5. size(): Devuelves cuántos elementos hay. O(1).

📝 PSEUDOCÓDIGO
──────────────────────────────────────────────────────────────

// Validar paréntesis balanceados usando un Stack
función esValido(cadena):
    pila = crear Stack
    PARA cada caracter EN cadena:
        SI caracter == '(':
            pila.push('(')
        SINO SI caracter == ')':
            SI pila.isEmpty(): return FALSO
            pila.pop()
    return pila.isEmpty()

💻 CÓDIGO C++
──────────────────────────────────────────────────────────────

#include <iostream>
#include <stack>    // STL Stack
#include <vector>
using namespace std;

// Implementación manual simple usando un vector
class MyStack {
private:
    vector<int> data;
public:
    void push(int x) { data.push_back(x); }
    void pop() { if(!data.empty()) data.pop_back(); }
    int top() { return data.back(); }
    bool isEmpty() { return data.empty(); }
};

int main() {
    // 1. Usando STL Stack
    stack<int> s;
    
    s.push(10);
    s.push(20);
    s.push(30);

    cout << "Tope actual: " << s.top() << endl; // 30

    s.pop(); // Quita el 30
    cout << "Nuevo tope: " << s.top() << endl;  // 20

    // 2. Recorrer (vaciando)
    while(!s.empty()) {
        cout << s.top() << " ";
        s.pop();
    }
    cout << endl;

    return 0;
}

⏱️ COMPLEJIDAD (Big-O)
──────────────────────────────────────────────────────────────

  Operación        │ Tiempo  │ Espacio
  ─────────────────┼─────────┼──────────────
  push(x)          │ O(1)    │ O(1)
  pop()            │ O(1)    │ O(1)
  peek() / top()   │ O(1)    │ O(1)
  isEmpty()        │ O(1)    │ O(1)
  Total (n datos)  │ -       │ O(n)
\end{lstlisting}

\entrevistas{
✅ Problema de Oro: "Valid Parentheses". Si ves paréntesis, corchetes 
   o llaves en un problema, ¡usa un stack!

💡 Recursión vs Stack: Cualquier problema recursivo se puede resolver
   usando un Stack de forma iterativa. A veces esto evita el Stack 
   Overflow.

⚠️ Stack Underflow: Intentar hacer `pop()` o `top()` en una pila vacía.
   SIEMPRE verifica `!s.empty()` antes de acceder al tope.

🔑 Patrón "Monotonic Stack": Un stack donde los elementos siempre están
   en orden (creciente o decreciente). Clave para problemas como 
   "Next Greater Element" o "Daily Temperatures".
}

\resumen{
• LIFO: Last In, First Out.
• Todas las operaciones básicas son O(1) (instantáneas).
• Usos: Historial del navegador (Undo/Redo), procesar expresiones
  matemáticas, DFS (Depth First Search).
• Muy eficiente pero limitado: solo ves el tope.
• La base de la recursión en cualquier lenguaje.
════════════════════════════════════════════════════════════════
}

\end{document}